% Test of noteoverflow=foot (didactic): the copy of a footnote that isn't printed
% says nothing in the log, and the one printed says what TeX would.  Each margin
% footnote is set twice under foot, in the margin and for the foot, and either may
% be printed.  On page 1 the margin is free: footnote A, with a token too wide for
% the margin or the text, and footnote D, with a line ended by \break, which is
% underfull only when justified at the foot's width, go in the margin.  On page 2 a
% tall note fills the margin: footnotes B (too wide) and C (\break) go to the foot.
% See test-noteoverflow-foot.sh, which builds it with foot (the default), with
% noteoverflow=break and with carry, and test-noteoverflow-foot.py hbox.
\documentclass[a4paper]{memoir}
\ifdefined\testoptions\else\def\testoptions{}\fi
\expandafter\PassOptionsToPackage\expandafter{\testoptions}{didactic}
\usepackage{didactic}
\newcommand\filler{Some words to fill up the rest of this line of the text and
  also quite a few of the next one, as the words of a paragraph do.}
\pagestyle{plain}
\begin{document}
\filler\footnote{Note A, with
TOKENATOOWIDEFORTHEMARGINORTHETEXTXXXXXXXXXXXXXXXXXXXXXXXXXXXXXXXXXXXXXXXXXXXXXXXXXXXX
in it.} \filler\ \filler

\filler\footnote{Note D, a short line\break and then the rest.} \filler\ \filler
\clearpage
\noindent x\marginpar{TALL\par\rule{1pt}{.97\textheight}} \filler\ \filler

\filler\footnote{Note B, with
TOKENBTOOWIDEFORTHEMARGINORTHETEXTXXXXXXXXXXXXXXXXXXXXXXXXXXXXXXXXXXXXXXXXXXXXXXXXXXXX
in it.} \filler\ \filler

\filler\footnote{Note C, a short line\break and then the rest.} \filler\ \filler
\end{document}
